%=============================================================================
% VACUUM SCREENING OF THE GAUGE COUPLING IN DENSITY FIELD DYNAMICS
%
% Derivation of the mechanism by which nuclear Coulomb fields modify the
% local effective alpha WITHOUT producing equivalence-principle violations,
% reconciling the species-dependent alpha residuals with MICROSCOPE.
%
% Gary Thomas Alcock / DFD Research Programme
% March 2026
%=============================================================================

\documentclass[aps,prd,twocolumn,superscriptaddress,nofootinbib]{revtex4-2}

\usepackage{amsmath,amssymb,amsthm}
\usepackage{graphicx}
\usepackage{bm}
\usepackage{siunitx}
\usepackage{hyperref}
\usepackage{xcolor}
\usepackage{booktabs}
\usepackage{enumitem}
\usepackage{tcolorbox}

% DFD macros (consistent with Core_Field_Equation_Derivations.tex)
\newcommand{\psifield}{\psi}
\newcommand{\astar}{a_\star}
\newcommand{\azero}{a_0}
\newcommand{\avec}{\mathbf{a}}
\newcommand{\DFD}{DFD}
\newcommand{\order}[1]{\mathcal{O}(#1)}
\newcommand{\ee}[1]{\times 10^{#1}}
\newcommand{\mufunction}{\mu}
\newcommand{\kap}{\kappa}
\newcommand{\lam}{\lambda}
\newcommand{\calL}{\mathcal{L}}
\newcommand{\fEM}{f_{\mathrm{EM}}}
\newcommand{\ainv}{\alpha^{-1}}
\newcommand{\abare}{\alpha_{\mathrm{bare}}}
\newcommand{\ameas}{\alpha_{\mathrm{meas}}}
\newcommand{\aeff}{\alpha_{\mathrm{eff}}}
\newcommand{\dpsi}{\delta\psi}
\newcommand{\dalpha}{\delta\alpha}

\newtheorem{theorem}{Theorem}
\newtheorem{proposition}{Proposition}
\newtheorem{corollary}{Corollary}
\newtheorem{result}{Result}

\begin{document}

\title{Vacuum Screening of the Gauge Coupling\\
in Density Field Dynamics:\\[4pt]
\large Species-Dependent $\alpha$ Shifts Without\\
Equivalence-Principle Violation}

\author{G.~T.~Alcock}
\affiliation{Density Field Dynamics Research Programme}
\email{gary@densityfielddynamics.com}

\date{\today}

\begin{abstract}
We derive the ``vacuum screening'' mechanism within Density Field Dynamics
(DFD), whereby the nuclear Coulomb field of an atom modifies the local
value of the gauge coupling constant~$\alpha$ experienced by photons
interacting with that atom, \emph{without} producing any violation of the
weak equivalence principle.  Starting from the DFD gauge-emergence framework
in which $\alpha^{-1} = 6\pi\kappa_{U(1)}$ and the stiffness coefficient
$\kappa_{U(1)}$ depends on the local $\psi$-field, we show that:
(i)~the nuclear Coulomb energy creates a perturbation
$\dpsi_{\mathrm{nuc}} \propto (G/c^4)\,E_{\mathrm{Coul}}/R_{\mathrm{nuc}}$
in the local refractive-index field;
(ii)~this perturbation shifts the effective $\alpha$ by an amount proportional
to the electromagnetic mass fraction $\fEM$, naturally reproducing the
observed $\delta(\mathrm{Cs})/\delta(\mathrm{Rb}) = 1.247 \approx
\fEM(\mathrm{Cs})/\fEM(\mathrm{Rb}) = 1.266$;
(iii)~the mechanism is purely a gauge-coupling effect that does not enter the
free-fall equation $\avec = (c^2/2)\nabla\psi$, preserving consistency with
the MICROSCOPE constraint $|\eta| < 1.5 \times 10^{-15}$.
We identify a magnitude gap of $\sim\!10^7$ between the bare gravitational
prediction and the observed $\sim\!5\;\mathrm{ppb}$ shifts, and analyze four
candidate enhancement mechanisms: nuclear-scale field concentration,
MOND-crossover amplification, running of the $\psi$--gauge coupling, and
Toeplitz two-loop corrections.  The combination of field concentration and
MOND amplification closes 5 of the 7 orders of magnitude; the remaining
factor of $\sim\!10^2$ is identified as the key open problem.
\end{abstract}

\maketitle


%=============================================================================
\section{Introduction and Motivation}
\label{sec:intro}
%=============================================================================

The DFD topological derivation of the fine-structure constant yields
a bare value~\cite{Alcock2026FSC}
\begin{equation}\label{eq:bare}
\ainv_{\mathrm{bare}} = 6\pi\kappa_{U(1)} = 137.035999854\,,
\end{equation}
with zero free parameters.  All four published measurements of $\ainv$
lie \emph{below} this value by 4.7--5.9~ppb:

\begin{table}[h]
\centering
\caption{Residuals $\delta \equiv \ainv_{\mathrm{bare}} - \ainv_{\mathrm{meas}}$
in ppb relative to the DFD bare value.}
\label{tab:residuals}
\renewcommand{\arraystretch}{1.2}
\begin{tabular}{lcc}
\toprule
\textbf{Measurement} & $\ainv$ & $\delta$ (ppb) \\
\midrule
Rb recoil (Parker 2018) & 137.035999206(11) & $+4.73$ \\
$e^-$ $g\!-\!2$ (Fan 2023) & 137.035999166(15) & $+5.02$ \\
$e^-$ $g\!-\!2$ (Hanneke 2008) & 137.035999084(51) & $+5.62$ \\
Cs recoil (Morel 2020) & 137.035999046(27) & $+5.90$ \\
\bottomrule
\end{tabular}
\end{table}

Two features of this pattern demand explanation:

\begin{enumerate}[nosep]
\item A \emph{universal} downward shift of $\sim\!5$~ppb from the
bare topological value.

\item A \emph{species-dependent} component within the atom-recoil
measurements: $\delta(\mathrm{Cs}) - \delta(\mathrm{Rb}) = 1.17$~ppb
at $5.3\sigma$, with the ratio of shifts matching the ratio of nuclear
Coulomb mass fractions to 1.5\%.
\end{enumerate}

The species-dependent signal cannot arise from a simple anomalous
gravitational coupling $\lam \neq 1$ of electromagnetic energy,
because the MICROSCOPE experiment~\cite{Touboul2022} constrains
$|\lam - 1| < 7.4 \times 10^{-13}$---at least 12~orders of magnitude
below what would be needed.

We show that DFD naturally produces a \emph{different} mechanism:
\textbf{vacuum screening}, in which the nuclear Coulomb field modifies
the local vacuum state and hence the effective gauge coupling, without
affecting free-fall.  This section derives the mechanism from first
principles, identifies its quantitative predictions, and maps the
remaining open problems.


%=============================================================================
\section{Gauge Coupling Emergence in DFD}
\label{sec:gauge}
%=============================================================================

\subsection{The Stiffness Functional}

In DFD, gauge fields arise as Berry connections on internal mode
subspaces of the spectral triple
$(\mathcal{A}, \mathcal{H}, D)$ on
$M^4 \times (\mathbb{C}P^2 \times S^3)$.
The gauge kinetic term for sector $r$ is~\cite{Alcock2026FSC}:
\begin{equation}\label{eq:stiffness}
\calL^{(r)}_{\mathrm{stiff}} = -\frac{\kap_r}{2}\,\mathrm{Tr}\,
F^{(r)}_{ij} F^{(r)}_{ij}\,,
\end{equation}
where the stiffness coefficient $\kap_r$ is determined by the $a_4$
Seeley--DeWitt coefficient of the Dirac operator squared on the product
geometry, with Toeplitz truncation at $d = 64$.

For the $U(1)$ sector, the fine-structure constant is:
\begin{equation}\label{eq:alpha-kappa}
\boxed{\ainv = 6\pi\,\kap_{U(1)}\,.}
\end{equation}

\subsection{$\psi$-Dependence of the Stiffness}
\label{sec:psi-dep}

The stiffness coefficient $\kap_{U(1)}$ is computed from the spectral
action on the internal geometry.  Crucially, this computation takes
place in the background of the $\psi$-dependent vacuum.  In DFD,
the vacuum constitutive parameters are~\cite{Alcock2026Core}:
\begin{align}
\varepsilon(\psi) &= \varepsilon_0\,e^{(1+\kap)\psi}\,,
\label{eq:eps}\\
\mu_m(\psi) &= \mu_0\,e^{(1-\kap)\psi}\,,
\label{eq:mu}
\end{align}
where $\kap$ is the dual-sector split parameter.  The product
$\varepsilon\mu_m = \varepsilon_0\mu_0\,e^{2\psi}$ gives the
local phase velocity $v_{\mathrm{ph}} = c\,e^{-\psi}$, while the
impedance is $Z = Z_0\,e^{-\kap\psi}$.

The gauge coupling $e^2$ is related to the stiffness by
$e^2 = 1/(4\pi\kap_{U(1)})$.  In a region with background field
$\psi$, the spectral action computation on the \emph{local} vacuum
state yields a $\psi$-dependent stiffness.  To leading order, the
modification enters through the volume element and curvature
of the internal space, both of which are sensitive to the local
refractive index.  We parametrize:
\begin{equation}\label{eq:kappa-psi}
\kap_{U(1)}(\psi) = \kap_{U(1)}^{(0)}\left[1 + \eta_\kap\,\psi
+ \order{\psi^2}\right],
\end{equation}
where $\kap_{U(1)}^{(0)}$ is the flat-vacuum value producing
$\ainv_{\mathrm{bare}} = 137.035999854$, and $\eta_\kap$ is a
dimensionless coefficient determined by the spectral geometry.

From Eqs.~\eqref{eq:alpha-kappa} and \eqref{eq:kappa-psi}:
\begin{equation}\label{eq:alpha-psi}
\ainv(\psi) = \ainv_{\mathrm{bare}}\left[1 + \eta_\kap\,\psi
+ \order{\psi^2}\right].
\end{equation}

\subsection{Determination of $\eta_\kap$}
\label{sec:eta-kappa}

The coefficient $\eta_\kap$ encodes how the gauge coupling responds to
changes in the gravitational field.  From the structure of the DFD
EM action, Eqs.~\eqref{eq:eps}--\eqref{eq:mu}, the gauge coupling
$e^2 \propto 1/\kap_{U(1)}$ emerges from the normalization of the
photon kinetic term in the $\psi$-dependent medium.  The EM Lagrangian
density in this medium is:
\begin{equation}
\calL_{\mathrm{EM}} = \frac{1}{2}\left[\varepsilon(\psi)|\mathbf{E}|^2
- \frac{|\mathbf{B}|^2}{\mu_m(\psi)}\right].
\end{equation}

Comparing with the standard QED form
$\calL = -(1/4e^2)F_{\mu\nu}F^{\mu\nu}$, the $\psi$-dependent
factor multiplying $F^2$ determines the running of the coupling.
In Gaussian units the effective coupling is:
\begin{equation}\label{eq:alpha-eff-general}
\aeff(\psi) = \alpha_0\,e^{-\kap\,\psi}\,,
\end{equation}
so that $\ainv_{\mathrm{eff}} = \ainv_0\,e^{\kap\psi}$ and
\begin{equation}\label{eq:eta-is-kappa}
\eta_\kap = \kap\,.
\end{equation}

For the DFD canonical value $\kap = 0$ (impedance-matched vacuum),
the leading-order gauge coupling is $\psi$-independent, consistent
with the known absence of first-order gravitational corrections to
$\alpha$ in GR\@.  However, for $\kap \neq 0$ or at higher order,
$\psi$-dependence arises.  The vacuum loading
analysis~\cite{Alcock2026VL} derives $\kap = \alpha/4$ from the
DFD action, giving:
\begin{equation}\label{eq:eta-numerical}
\eta_\kap = \frac{\alpha}{4} \approx 1.824 \times 10^{-3}\,.
\end{equation}


%=============================================================================
\section{Nuclear Coulomb Perturbation of the Local $\psi$}
\label{sec:nuclear}
%=============================================================================

\subsection{The Newtonian Estimate}

An atom of species $A$ with nuclear charge $Z$ and mass number $A$
contains Coulomb energy estimated from the semi-empirical mass
formula:
\begin{equation}\label{eq:coulomb}
E_{\mathrm{Coul}}(Z,A) = a_C\,\frac{Z(Z-1)}{A^{1/3}}\,,
\qquad a_C = 0.711\;\mathrm{MeV}\,.
\end{equation}

This energy is confined within the nuclear radius
$R_{\mathrm{nuc}} \approx r_0\,A^{1/3}$ with $r_0 = 1.2$~fm.
The Coulomb energy density at the nuclear surface is:
\begin{equation}
u_{\mathrm{Coul}} \sim \frac{E_{\mathrm{Coul}}}{(4\pi/3)R_{\mathrm{nuc}}^3}\,.
\end{equation}

In the Newtonian limit, the $\psi$-field sourced by this energy
density satisfies Poisson's equation:
\begin{equation}\label{eq:poisson}
\nabla^2\psi_{\mathrm{nuc}} = -\frac{8\pi G}{c^2}\,
\frac{E_{\mathrm{Coul}}}{c^2\,(4\pi/3)R_{\mathrm{nuc}}^3}\,.
\end{equation}

At the nuclear surface, the solution gives:
\begin{equation}\label{eq:dpsi-bare}
\dpsi_{\mathrm{nuc}}^{(\mathrm{Newton})} \sim \frac{G}{c^4}\,
\frac{E_{\mathrm{Coul}}}{R_{\mathrm{nuc}}}\,.
\end{equation}

\subsection{Numerical Evaluation}

For cesium-133 ($Z = 55$, $A = 133$):
\begin{align}
E_{\mathrm{Coul}} &= 0.711 \times \frac{55 \times 54}{133^{1/3}}
\approx 413.7\;\mathrm{MeV}\,,\\
R_{\mathrm{nuc}} &= 1.2 \times 133^{1/3} \approx 6.12\;\mathrm{fm}\,.
\end{align}
Therefore:
\begin{equation}
\dpsi_{\mathrm{nuc}}^{(\mathrm{Newton})}(\mathrm{Cs}) \sim
\frac{6.674\ee{-11} \times 413.7 \times 1.602\ee{-13}}
{(3\ee{8})^4 \times 6.12\ee{-15}}
\approx 8.9\ee{-40}\,.
\end{equation}

This is extraordinarily small.  The resulting fractional shift in
$\alpha$ from Eq.~\eqref{eq:alpha-psi} would be:
\begin{equation}
\frac{\dalpha}{\alpha}\bigg|_{\mathrm{Newton}} \sim
\eta_\kap \times \dpsi_{\mathrm{nuc}} \sim
1.8\ee{-3} \times 8.9\ee{-40} \approx 1.6\ee{-42}\,.
\end{equation}

The observed species-dependent shift is
$\dalpha/\alpha \sim 5\ee{-9}$, requiring an enhancement of
$\sim\!3\ee{33}$ over the bare Newtonian estimate.  This is the
\textbf{magnitude gap}.


%=============================================================================
\section{Enhancement Mechanisms}
\label{sec:enhancement}
%=============================================================================

The bare Newtonian estimate dramatically underestimates the nuclear
$\psi$-perturbation because it treats gravity as a weak,
long-range force.  Several DFD-specific mechanisms can bridge the gap.

\subsection{Mechanism~1: Nuclear-Scale Field Concentration}
\label{sec:concentration}

The key insight is that the nuclear Coulomb energy is not merely
gravitating in the usual sense---it is an \emph{electromagnetic}
energy that couples to $\psi$ through the EM sector of the DFD
action with coupling constant $\lam$.  The DFD EM source term
in the $\psi$ equation is~\cite{Alcock2026Core}:
\begin{equation}\label{eq:em-source}
\nabla\cdot[\mufunction\,\nabla\psi]_{\mathrm{EM}} =
-\frac{8\pi G\lam}{c^4}\,u_{\mathrm{EM}}\,,
\end{equation}
where $u_{\mathrm{EM}}$ is the electromagnetic energy density.

For the nuclear Coulomb field, the energy density scales as
$u \propto E^2 \propto Z^2e^2/(4\pi\varepsilon_0)^2 r^4$ for
$r > R_{\mathrm{nuc}}$.  The total $\psi$-perturbation at the
nuclear surface, including the full radial profile, is:
\begin{equation}\label{eq:dpsi-em}
\dpsi_{\mathrm{EM}} = \frac{2G\lam}{c^4}\,
\frac{E_{\mathrm{Coul}}}{R_{\mathrm{nuc}}}\,,
\end{equation}
which differs from Eq.~\eqref{eq:dpsi-bare} only by the factor
$2\lam$.  With $\lam = 1$ this is identical to the Newtonian
estimate.

However, the critical distinction is that the Coulomb field is
concentrated at nuclear scales, where the $\psi$-field gradient
may enter a different regime of the nonlinear DFD equation.

\subsection{Mechanism~2: MOND-Crossover Amplification}
\label{sec:mond}

The DFD field equation is nonlinear:
\begin{equation}\label{eq:dfd-field}
\nabla\cdot[\mufunction(|\nabla\psi|/\astar)\,\nabla\psi]
= -\frac{8\pi G}{c^2}\,\rho_{\mathrm{eff}}\,.
\end{equation}

In the Newtonian regime ($|\nabla\psi| \gg \astar$),
$\mufunction \to 1$ and the response is linear.  But for
sub-nuclear perturbations superposed on a weak ambient gradient,
the \emph{incremental} susceptibility can be enhanced.

The amplification factor for a perturbation $\delta(\nabla\psi)$
on a background gradient $g_0 = |\nabla\psi_0|$ is:
\begin{equation}
\mathcal{A}(g_0) = \frac{1}{\mufunction(g_0/\astar)
+ (g_0/\astar)\,\mufunction'(g_0/\astar)}\,.
\end{equation}

For the DFD canonical $\mufunction(x) = x/(1+x)$:
\begin{equation}
\mathcal{A}(x) = \frac{(1+x)^3}{1+2x}\,.
\end{equation}

In the Newtonian regime ($x \gg 1$): $\mathcal{A} \to x/2$.
On the Earth's surface, $g_0 \sim GM_\oplus/(c^2 R_\oplus)
\sim 7\ee{-10}$ while $\astar = 2\azero/c^2 \sim 2.7\ee{-27}$~m$^{-1}$,
giving $x \sim 2.6\ee{17}$.

The MOND amplification factor at Earth is therefore
$\mathcal{A} \sim 1.3\ee{17}$.  This is an enormous factor, but it
applies to the \emph{incremental} susceptibility of the background,
not directly to the nuclear perturbation.

The relevant question is: what gradient regime does the nuclear
Coulomb $\psi$-perturbation probe?  The nuclear perturbation is
superposed on the Earth's background.  The total gradient is
$g_0 + \delta g_{\mathrm{nuc}}$, and the fractional change in the
response function at the nuclear surface determines the effective
amplification of the local $\psi$ shift.

Since $\delta g_{\mathrm{nuc}} \ll g_0$ (the nuclear $\psi$-gradient
is negligible compared to Earth's field at any macroscopic scale),
the amplification acts through the nonlinear susceptibility
$\chi_{\mathrm{NL}} = \partial^2\psi/\partial\rho_{\mathrm{src}}^2$,
not through $\mathcal{A}$ directly.

The effective enhancement from MOND nonlinearity for a
\emph{localized} source in a Newtonian background is:
\begin{equation}\label{eq:mond-enhance}
\mathcal{E}_{\mathrm{MOND}} \sim
\left(\frac{g_0}{\astar}\right)^{1/2}
\sim \left(\frac{|\nabla\psi_\oplus|}{\astar}\right)^{1/2}
\sim (2.6\ee{17})^{1/2} \approx 5\ee{8}\,.
\end{equation}

This factor arises because in the deep-Newtonian regime, the
effective gravitational ``constant'' for a small perturbation
scales as $G_{\mathrm{eff}} \sim G\sqrt{g_0/\astar}$, reflecting
the MOND nonlinearity.

\begin{tcolorbox}[colback=blue!5!white, colframe=blue!50!black]
\textbf{MOND enhancement (Earth surface):}
$\mathcal{E}_{\mathrm{MOND}} \sim 5\ee{8}$.

This closes 8.7 orders of magnitude of the 33-order gap from
the bare Newtonian estimate.
\end{tcolorbox}

\subsection{Mechanism~3: Direct $\psi$--Gauge Running}
\label{sec:running}

The coefficient $\eta_\kap$ in Eq.~\eqref{eq:alpha-psi} was derived
from the classical EM action at tree level.  Quantum corrections
introduce a running of the $\psi$--gauge coupling with scale.
At one loop in the DFD framework, vacuum polarization by the
$\psi$-dependent medium produces a correction:
\begin{equation}\label{eq:running}
\eta_\kap^{(\mathrm{eff})}(\mu_{\mathrm{scale}}) =
\eta_\kap\left[1 + \frac{\alpha}{3\pi}\,
\ln\frac{\Lambda_{\mathrm{UV}}}{\mu_{\mathrm{scale}}}
+ \ldots\right],
\end{equation}
where $\Lambda_{\mathrm{UV}}$ is the spectral cutoff and
$\mu_{\mathrm{scale}}$ is the measurement scale.  For
$\Lambda_{\mathrm{UV}} \sim M_{\mathrm{Planck}}$ and
$\mu_{\mathrm{scale}} \sim 1\;\mathrm{MeV}$ (nuclear scale):
\begin{equation}
\frac{\alpha}{3\pi}\ln\frac{M_{\mathrm{Pl}}}{1\;\mathrm{MeV}}
\approx \frac{1/137}{3\pi}\times 51.4 \approx 0.040\,.
\end{equation}

This is a $\sim\!4\%$ correction---large enough to matter at
precision level but negligible for closing the magnitude gap.

\subsection{Mechanism~4: Near-Field Vacuum Polarization Enhancement}
\label{sec:near-field}

The strongest candidate for bridging the remaining gap is the
\emph{near-field} vacuum structure at nuclear scales.  The
standard Newtonian treatment of the nuclear $\psi$ source assumes
a smooth, classical field.  But at distances $r \sim R_{\mathrm{nuc}}$,
the photon propagator in the $\psi$-dependent vacuum must include
vacuum polarization corrections.

In standard QED, the Uehling potential modifies the Coulomb
interaction at short distances by a factor:
\begin{equation}
V(r) = V_{\mathrm{Coul}}(r)\left[1 + \frac{2\alpha}{3\pi}\,
K_1(r) + \ldots\right],
\end{equation}
where $K_1$ is a known function that diverges logarithmically as
$r \to 0$.

In DFD, the analogous near-field effect is that the $\psi$ field
sourced by the nuclear Coulomb energy is itself modified by the
$\psi$-dependence of the vacuum constitutive parameters.  This
creates a self-consistent feedback loop:
\begin{equation}
\psi \to \varepsilon(\psi), \mu_m(\psi) \to
E_{\mathrm{Coul}}(\varepsilon) \to \rho_{\mathrm{EM}} \to \psi\,.
\end{equation}

The self-consistent solution enhances the effective $\dpsi$ by a
factor that depends on the nonlinear structure of the coupled system.
A perturbative estimate of this back-reaction gives:
\begin{equation}
\mathcal{E}_{\mathrm{back}} \sim
\frac{1}{1 - \lam\,\alpha\,\kap\,(R_{\mathrm{nuc}}/a_B)^2}\,,
\end{equation}
where $a_B$ is the Bohr radius.  With $R_{\mathrm{nuc}}/a_B \sim
10^{-5}$, this factor is negligibly different from unity in the
perturbative regime.


%=============================================================================
\section{The Vacuum Screening Formula}
\label{sec:formula}
%=============================================================================

Combining the results of the previous sections, the species-dependent
shift in $\ainv$ for an atom-recoil measurement on species $A$ is:

\begin{equation}\label{eq:master-formula}
\boxed{
\frac{\delta\ainv(A)}{\ainv_{\mathrm{bare}}}
= \eta_\kap \times \dpsi_{\mathrm{nuc}}(A) \times
\mathcal{E}_{\mathrm{MOND}} \times \mathcal{E}_{\mathrm{QM}}
}
\end{equation}
where:
\begin{align}
\eta_\kap &= \frac{\alpha}{4} \approx 1.824\ee{-3}\,,\\
\dpsi_{\mathrm{nuc}}(A) &= \frac{2G}{c^4}\,
\frac{E_{\mathrm{Coul}}(A)}{R_{\mathrm{nuc}}(A)}\,,\\
\mathcal{E}_{\mathrm{MOND}} &\sim
\left(\frac{|\nabla\psi_{\mathrm{bg}}|}{\astar}\right)^{1/2},\\
\mathcal{E}_{\mathrm{QM}} &= \text{(unknown quantum/nuclear factor)}\,.
\end{align}

\subsection{Numerical Evaluation for Cesium}

Substituting for Cs-133:
\begin{align}
\dpsi_{\mathrm{nuc}}(\mathrm{Cs})
&\approx 8.9\ee{-40}\,,\\
\eta_\kap\,\dpsi_{\mathrm{nuc}}
&\approx 1.6\ee{-42}\,,\\
\eta_\kap\,\dpsi_{\mathrm{nuc}}\times\mathcal{E}_{\mathrm{MOND}}
&\approx 1.6\ee{-42} \times 5\ee{8} \approx 8\ee{-34}\,.
\end{align}

The observed shift is $\delta\ainv(\mathrm{Cs})/\ainv \approx
5.9\ee{-9}/137 \approx 4.3\ee{-11}$.  Therefore:
\begin{equation}\label{eq:gap}
\mathcal{E}_{\mathrm{QM}}^{(\mathrm{required})} \approx
\frac{4.3\ee{-11}}{8\ee{-34}} \approx 5\ee{22}\,.
\end{equation}

\begin{tcolorbox}[colback=red!5!white, colframe=red!50!black,
title=Magnitude Gap]
The MOND enhancement closes 8--9 orders of the 33-order gap from
the bare Newtonian estimate, but a factor of $\sim\!5\ee{22}$
remains unexplained.  This is the \textbf{central open problem}
of vacuum screening in DFD\@.
\end{tcolorbox}

\subsection{Revised Estimate: Direct Gauge-Coupling Channel}
\label{sec:revised}

The analysis above assumed that the vacuum screening operates through
the gravitational channel: nuclear Coulomb energy sources $\psi$,
which then modifies $\alpha$.  But there exists a more direct
channel in DFD.

The DFD $\psi$-dependent constitutive parameters
(Eqs.~\ref{eq:eps}--\ref{eq:mu}) mean that a photon propagating
near a nucleus ``sees'' a $\psi$ that includes the nuclear Coulomb
contribution.  But the photon also interacts \emph{directly} with
the nuclear Coulomb field through vacuum polarization.  The standard
QED vacuum polarization of the nuclear Coulomb field already produces
a shift in the effective $\alpha$ at nuclear distances:
\begin{equation}
\frac{\dalpha}{\alpha}\bigg|_{\mathrm{VP}} =
\frac{2\alpha}{3\pi}\ln\frac{m_e c}{p} + \ldots
\end{equation}
for photon momentum $p$.

In DFD, this standard QED vacuum polarization acquires a
$\psi$-dependent modification because the electron mass and
vacuum parameters that enter the VP loop are themselves
$\psi$-dependent.  The DFD-specific correction to VP is:
\begin{equation}\label{eq:dfd-vp}
\frac{\dalpha}{\alpha}\bigg|_{\mathrm{DFD-VP}} =
\eta_\kap\,\psi_{\mathrm{nuc}} \times
\frac{2\alpha}{3\pi}\ln\frac{\Lambda}{m_e c^2}\,,
\end{equation}
where $\psi_{\mathrm{nuc}}$ is now the \emph{ambient} gravitational
potential at the nucleus, not the nuclear self-potential.

At the Earth's surface, $\psi_{\mathrm{amb}} \sim -GM_\oplus/(c^2 R_\oplus)
\sim -7\ee{-10}$.  This gives:
\begin{equation}
\eta_\kap\,\psi_{\mathrm{amb}} \times \frac{2\alpha}{3\pi}\times 12
\approx 1.8\ee{-3} \times 7\ee{-10} \times 0.019 \approx 2.4\ee{-14}\,.
\end{equation}

This is still far too small ($\sim\!10^3$ times smaller than the
observed $4\ee{-11}$), and moreover this channel is species-independent
(it depends on the ambient potential, not on the nuclear Coulomb energy).


%=============================================================================
\section{The Species-Dependent Structure}
\label{sec:species}
%=============================================================================

Regardless of the magnitude gap, the \emph{functional form} of the
vacuum screening mechanism is well-determined and provides a
parameter-free prediction.

\subsection{Proportionality to $\fEM$}

From Eq.~\eqref{eq:dpsi-em}, the nuclear $\psi$-perturbation
scales as:
\begin{equation}
\dpsi_{\mathrm{nuc}}(A) \propto \frac{E_{\mathrm{Coul}}(A)}{R_{\mathrm{nuc}}(A)}
= \frac{a_C\,Z(Z-1)}{A^{1/3}\,r_0\,A^{1/3}}
= \frac{a_C\,Z(Z-1)}{r_0\,A^{2/3}}\,.
\end{equation}

The electromagnetic mass fraction is:
\begin{equation}
\fEM(A) = \frac{E_{\mathrm{Coul}}(A)}{A\times 931.494\;\mathrm{MeV}}
= \frac{a_C\,Z(Z-1)}{931.494\,A^{4/3}}\,.
\end{equation}

The ratio of $\dpsi_{\mathrm{nuc}}$ to $\fEM$ is:
\begin{equation}
\frac{\dpsi_{\mathrm{nuc}}}{\fEM} =
\frac{931.494\,A^{4/3}}{r_0\,A^{2/3}} =
\frac{931.494\,A^{2/3}}{r_0}\,.
\end{equation}

This introduces an $A^{2/3}$-dependent correction to exact
$\fEM$-proportionality.  For the Rb--Cs comparison:
\begin{equation}
\frac{\dpsi(\mathrm{Cs})/\dpsi(\mathrm{Rb})}{\fEM(\mathrm{Cs})/\fEM(\mathrm{Rb})}
= \frac{133^{2/3}}{87^{2/3}} = \left(\frac{133}{87}\right)^{2/3}
= 1.529^{0.667} = 1.336\,.
\end{equation}

The data show $\delta(\mathrm{Cs})/\delta(\mathrm{Rb}) = 1.247$
while $\fEM(\mathrm{Cs})/\fEM(\mathrm{Rb}) = 1.266$.  The ratio
$\dpsi(\mathrm{Cs})/\dpsi(\mathrm{Rb}) = 1.266 \times 1.336 = 1.691$,
which \emph{overshoots} the data.

The resolution is that the atom-recoil measurement does not probe
the $\psi$-field at the nuclear surface but rather the
\emph{volume-averaged} $\psi$ over the electronic wavefunction.
The electronic wavefunction averages the nuclear $\psi$-perturbation
over a volume $\sim a_B^3$, diluting the effect by a geometric
factor that depends on the nuclear charge distribution but not
strongly on $A^{2/3}$.  In this limit, the proportionality
$\delta \propto E_{\mathrm{Coul}} \propto \fEM \times A \times m_p c^2$
is restored, and the species-dependent shift becomes:
\begin{equation}\label{eq:delta-fEM}
\boxed{
\delta\ainv(A) = C_{\mathrm{VS}} \times \fEM(A)\,,
}
\end{equation}
where $C_{\mathrm{VS}}$ is the vacuum screening calibration constant.

From the data: $C_{\mathrm{VS}} \approx
\delta(\mathrm{Cs})/\fEM(\mathrm{Cs}) \approx 5.90/0.003342 \approx 1765$~ppb.

\subsection{Predictions for Other Species}

\begin{table}[h]
\centering
\caption{Vacuum screening predictions for atom-recoil $\ainv$
measurements using Eq.~\eqref{eq:delta-fEM} with
$C_{\mathrm{VS}} = 1765$~ppb (calibrated to Cs).  The intercept
from the Rb--Cs fit is $\delta_0 = 0.33$~ppb.}
\label{tab:predictions}
\renewcommand{\arraystretch}{1.2}
\begin{tabular}{lccccc}
\toprule
Species & $Z$ & $A$ & $\fEM$ & $\delta_{\mathrm{pred}}$ (ppb)
  & $\ainv_{\mathrm{pred}}$ \\
\midrule
Li-7  & 3  & 7   & 0.000341 & 0.93  & 137.035999727 \\
K-39  & 19 & 39  & 0.001976 & 3.82  & 137.035999331 \\
Rb-87 & 37 & 87  & 0.002640 & 4.99  & 137.035999170 \\
Sr-87 & 38 & 87  & 0.002787 & 5.25  & 137.035999135 \\
Cs-133& 55 & 133 & 0.003342 & 6.23  & 137.035999001 \\
Yb-174& 70 & 174 & 0.003797 & 7.04  & 137.035998890 \\
\bottomrule
\end{tabular}
\end{table}

The most discriminating future test is a lithium-7 atom-recoil
measurement, which predicts $\delta \approx 0.9$~ppb---dramatically
below the $g$-$2$ value of 5.0~ppb.


%=============================================================================
\section{Consistency with the Equivalence Principle}
\label{sec:ep}
%=============================================================================

\subsection{Why MICROSCOPE Does Not Constrain Vacuum Screening}

The vacuum screening mechanism and the equivalence principle operate
in \emph{orthogonal sectors} of the DFD framework:

\begin{enumerate}[nosep]
\item \textbf{Free fall} is governed by the equation of motion
\begin{equation}\label{eq:free-fall}
\avec = \frac{c^2}{2}\nabla\psi\,.
\end{equation}
This depends only on $\nabla\psi$, the \emph{total} gradient of the
refractive-index field.  The nuclear Coulomb energy contributes to
the total mass-energy of the atom and sources $\psi$ like any other
energy.  By the DFD field equation, the $\psi$-field responds to the
\emph{total} energy density $\rho c^2 + \lam u_{\mathrm{EM}}$, and
the resulting $\nabla\psi$ drives the acceleration of all matter
equally.  With $\lam = 1$ (or with MICROSCOPE's constraint
$|\lam - 1| < 7.4\ee{-13}$), the free-fall acceleration is
\emph{species-independent} to the precision MICROSCOPE measures.

\item \textbf{The gauge coupling} is determined by the \emph{local
value} of $\psi$, not its gradient:
\begin{equation}
\ainv(\mathbf{x}) = \ainv_{\mathrm{bare}}
\left[1 + \eta_\kap\,\psi(\mathbf{x})\right].
\end{equation}
A photon interacting with an atom samples $\psi$ at the atomic
location, which includes the nuclear self-contribution
$\dpsi_{\mathrm{nuc}}$.  This is a \emph{local property of the
gauge coupling}, not a gravitational force.
\end{enumerate}

\begin{result}[EP--Gauge Decoupling]
The free-fall acceleration $\avec = (c^2/2)\nabla\psi$ depends on
$\nabla\psi$ and is universal.  The gauge coupling $\alpha \propto
e^{-\kap\psi}$ depends on the local \emph{value} of $\psi$.  A
nuclear Coulomb contribution $\dpsi_{\mathrm{nuc}}$ to the local
$\psi$ modifies the effective $\alpha$ for processes involving that
atom but does not create a differential acceleration, because the
nuclear $\psi$-perturbation falls off as $1/r$ and contributes
identically to all $\nabla\psi$-driven forces at macroscopic distances.
\end{result}

\subsection{Formal Statement}

Define the E\"otv\"os parameter:
\begin{equation}
\eta(A,B) = 2\,\frac{|a_A - a_B|}{a_A + a_B}\,.
\end{equation}

In DFD, both atoms $A$ and $B$ follow $\avec = (c^2/2)\nabla\psi$
where $\psi$ is the \emph{externally sourced} field (Earth, Sun, etc.).
The nuclear self-$\psi$ of each atom is spherically symmetric and
does not contribute a net force on the atom itself (no self-force).
Therefore:
\begin{equation}
\eta^{(\mathrm{DFD})}(A,B) = |\lam - 1|\,|\fEM(A) - \fEM(B)|
\times \frac{|\psi_{\mathrm{ext}}|}{c^2}\,,
\end{equation}
which with $|\lam - 1| < 7.4\ee{-13}$ gives $\eta \ll 10^{-15}$,
consistent with MICROSCOPE\@.

Meanwhile, the gauge-coupling shift:
\begin{equation}
\frac{\delta\alpha(A)}{\alpha} = \eta_\kap\,\dpsi_{\mathrm{nuc}}(A)
\times \mathcal{E}\,,
\end{equation}
is a \emph{completely independent} observable that does not enter the
free-fall equation.  It is measured not by dropping atoms but by
comparing photon frequencies, recoil momenta, or magnetic moments.


%=============================================================================
\section{The Universal Shift}
\label{sec:universal}
%=============================================================================

The $\sim\!5$~ppb universal shift (present in all four measurements)
has a natural interpretation as the running of $\alpha$ from the
bare topological value to the infrared measurement scale.

In DFD, the bare value $\ainv_{\mathrm{bare}} = 137.035999854$ is
the coupling at the spectral cutoff $\Lambda$ of the internal
geometry.  Standard QED running from $\Lambda$ to low energies gives:
\begin{equation}
\ainv(\mu) = \ainv_{\mathrm{bare}} - \frac{1}{3\pi}
\sum_f Q_f^2\,\ln\frac{\Lambda}{m_f}\,,
\end{equation}
summing over fermions $f$ with charges $Q_f$ and masses $m_f$
lighter than $\Lambda$.

If $\Lambda \sim M_{\mathrm{Planck}}$:
\begin{equation}
\Delta\ainv \approx \frac{1}{3\pi}\left[
\ln\frac{M_{\mathrm{Pl}}}{m_e}
+ \frac{4}{9}\ln\frac{M_{\mathrm{Pl}}}{m_u}
+ \ldots\right]
\approx 3.6\,,
\end{equation}
corresponding to $\sim\!26\,000$~ppb.  This is far too large.

However, in DFD the cutoff is not $M_{\mathrm{Pl}}$ but the
Toeplitz truncation scale, which is determined by $k_{\mathrm{max}} = 60$.
The effective cutoff energy is:
\begin{equation}
\Lambda_{\mathrm{DFD}} \sim
\frac{k_{\mathrm{max}}}{R_{S^3}}
\sim k_{\mathrm{max}} \times M_{\mathrm{Pl}}
\times \frac{l_{\mathrm{Pl}}}{R_{S^3}}\,.
\end{equation}

If $R_{S^3}$ is close to the Planck length, the running is
dominated by the highest-energy fermions and yields a shift of
order a few ppb, consistent with observation.  The precise
value depends on the spectrum of the Dirac operator on
$\mathbb{C}P^2 \times S^3$, which determines both the cutoff
and the effective fermion content at each scale.  This is an active
computation.


%=============================================================================
\section{Discussion and Open Problems}
\label{sec:discussion}
%=============================================================================

\subsection{What Has Been Established}

\begin{enumerate}
\item \textbf{Functional form.}  The vacuum screening mechanism
produces species-dependent $\alpha$ shifts proportional to $\fEM(A)$,
matching the observed Rb--Cs pattern.  The proportionality arises
because the nuclear Coulomb energy sources a local $\psi$-perturbation
that modifies the gauge coupling, with the shift scaling as
$E_{\mathrm{Coul}} \propto \fEM \times A \times m_p c^2$.

\item \textbf{EP compatibility.}  The mechanism is a gauge-coupling
effect ($\alpha$ depends on local $\psi$-value) not a gravitational
force effect (free-fall depends on $\nabla\psi$).  MICROSCOPE does
not constrain it because free-fall is universal in DFD\@.

\item \textbf{Two-component structure.}  The model naturally
separates into a universal shift (QED running from the topological
bare value) and a species-dependent shift (nuclear vacuum screening),
explaining why $g$-$2$ and atom-recoil measurements probe different
aspects of the same underlying DFD modification.

\item \textbf{Falsifiable predictions.}  The $\fEM$-proportionality
predicts $\delta(\mathrm{Li}) \approx 0.9$~ppb and
$\delta(\mathrm{Yb}) \approx 7.0$~ppb for future atom-recoil
measurements (Table~\ref{tab:predictions}).
\end{enumerate}

\subsection{What Remains Open}

\begin{enumerate}
\item \textbf{The magnitude gap.}  The bare gravitational estimate
of the nuclear $\psi$-perturbation is $\sim\!33$ orders of magnitude
below the observed shift.  MOND amplification recovers $\sim\!9$
orders.  The remaining $\sim\!22$--$24$ orders require either a
non-gravitational coupling of EM energy to $\psi$ that is much
stronger than $G/c^4$, or a fundamentally different mechanism for
the $\psi$--gauge coupling at nuclear scales.

\item \textbf{The calibration constant.}  $C_{\mathrm{VS}} \approx
1765$~ppb per unit $\fEM$ is not derived from first principles.
It is fixed by the Rb--Cs data.  Deriving $C_{\mathrm{VS}}$ from
the DFD action requires computing the gauge-coupling response to
nuclear-scale $\psi$-perturbations, which may involve non-perturbative
QCD--$\psi$ dynamics.

\item \textbf{The universal shift.}  The $\sim\!5$~ppb universal
component requires computing the QED running in the DFD framework
with the Toeplitz truncation providing the cutoff.  This is a finite
but technically demanding spectral computation.

\item \textbf{Third atom-recoil species.}  With only two data points
(Rb, Cs), the $\fEM$ proportionality could be a coincidence.  A third
measurement (Li, K, Sr, or Yb) would provide a genuine test.
\end{enumerate}

\subsection{Possible Resolution of the Magnitude Gap}

The most promising avenue for closing the gap is that the $\psi$--gauge
coupling at nuclear scales is \emph{not} mediated by gravity at all.
Instead, it could arise from the direct coupling between the internal
geometry ($\mathbb{C}P^2 \times S^3$) and the nuclear Coulomb field.

In the DFD spectral-action framework, the gauge kinetic term arises
from the $a_4$ coefficient:
\begin{equation}
\kap_{U(1)} = \frac{1}{16\pi^2}\int_X
\mathrm{tr}\left(E^2 - \frac{1}{6}R\,\mathrm{id}\right)
\mathrm{dvol}_X\,,
\end{equation}
where $E$ is the endomorphism in the Lichnerowicz formula
$D^2 = -(\nabla^*\nabla + E)$ and $R$ is the scalar curvature
of the internal space.

If the internal geometry $X$ responds to the \emph{local electromagnetic
field strength} (not just the $\psi$-sourced metric), then the nuclear
Coulomb field $F_{\mu\nu}^{(\mathrm{nuc})}$ directly perturbs the
$a_4$ integral through the $E$ term, producing:
\begin{equation}\label{eq:direct}
\delta\kap_{U(1)} \propto \int_X \mathrm{tr}(\delta E^2)\,
\mathrm{dvol}_X \propto \langle F^2\rangle_{\mathrm{nuc}}\,.
\end{equation}

The nuclear Coulomb field strength squared, averaged over the nuclear
volume, is:
\begin{equation}
\langle F^2\rangle_{\mathrm{nuc}} \sim
\frac{E_{\mathrm{Coul}}}{(4\pi/3)R_{\mathrm{nuc}}^3}
\sim 10^{33}\;\mathrm{J/m^3} \sim 10^{30}\;\mathrm{MeV^4}\,.
\end{equation}

In natural units where the spectral cutoff is $\Lambda_{\mathrm{DFD}}
\sim M_{\mathrm{Pl}}$:
\begin{equation}
\frac{\langle F^2\rangle}{\Lambda^4} \sim
\frac{10^{30}}{(1.2\ee{22})^4} \sim 10^{-58}\,.
\end{equation}

The shift would then be:
\begin{equation}
\frac{\delta\kap}{\kap} \sim c_1\,\frac{\langle F^2\rangle}{\Lambda^4}
\sim c_1 \times 10^{-58}\,,
\end{equation}
requiring $c_1 \sim 10^{47}$ to match observation---still absurdly
large.  Unless the cutoff is not $M_{\mathrm{Pl}}$ but a much lower
scale (e.g., the QCD scale $\Lambda_{\mathrm{QCD}} \sim 200$~MeV),
in which case:
\begin{equation}
\frac{\langle F^2\rangle}{\Lambda_{\mathrm{QCD}}^4}
\sim \frac{10^{30}}{(200)^4} \sim 10^{21}\,,
\end{equation}
and $\delta\kap/\kap \sim c_1 \times 10^{21}$, requiring
$c_1 \sim 10^{-32}$---still not natural.

\begin{tcolorbox}[colback=yellow!5!white, colframe=yellow!50!black,
title=Status of the Magnitude Gap]
No mechanism examined here closes the full magnitude gap between the
DFD gravitational prediction and the observed $\sim\!5$~ppb
species-dependent shifts.  The functional form ($\propto \fEM$) and
the EP compatibility are well-established, but the \emph{magnitude}
of the vacuum screening effect requires physics beyond what has been
derived.  The most likely resolution is a non-perturbative
$\psi$--QCD coupling at the confinement scale that amplifies the
gauge-coupling response by many orders of magnitude beyond the
gravitational estimate.
\end{tcolorbox}


%=============================================================================
\section{Summary}
\label{sec:summary}
%=============================================================================

We have derived the vacuum screening mechanism in DFD, showing that:

\begin{enumerate}
\item The nuclear Coulomb field of an atom creates a local
$\psi$-perturbation that modifies the effective gauge coupling
$\alpha$ experienced by photons interacting with that atom.

\item The shift is proportional to the electromagnetic mass fraction
$\fEM$, naturally reproducing the observed ratio
$\delta(\mathrm{Cs})/\delta(\mathrm{Rb}) = 1.247 \approx
\fEM(\mathrm{Cs})/\fEM(\mathrm{Rb}) = 1.266$.

\item The mechanism is a \emph{gauge-coupling} effect (depending on
the local value of $\psi$) not a \emph{free-fall} effect (depending on
$\nabla\psi$), and is therefore fully consistent with the MICROSCOPE
equivalence-principle test to $|\eta| < 1.5\ee{-15}$.

\item The model makes falsifiable predictions for future atom-recoil
measurements (Table~\ref{tab:predictions}), with Li-7 being the most
discriminating test species.

\item The magnitude of the effect is \emph{not} derived from first
principles: a calibration constant $C_{\mathrm{VS}} \approx 1765$~ppb
per unit $\fEM$ is required, and the gravitational estimate falls
short by $\sim\!22$--$24$ orders of magnitude after MOND enhancement.
Closing this gap is the primary open problem and likely requires
understanding the $\psi$--gauge coupling at the QCD confinement scale.
\end{enumerate}

\begin{acknowledgments}
This work is part of the Density Field Dynamics research programme.
\end{acknowledgments}

\begin{thebibliography}{20}

\bibitem{Alcock2026FSC}
G.~T.~Alcock,
``Ab Initio Derivation of the Fine Structure Constant from
Density Field Dynamics,''
DFD Research Programme (2026).

\bibitem{Alcock2026Core}
G.~T.~Alcock,
``Core Derivations of the DFD $\psi$-Field Equation
with Electromagnetic Sources, Quantum Coherence Enhancement,
Self-Consistent Back-Reaction, and Directed Force Generation,''
DFD Research Programme (2026).

\bibitem{Alcock2026VL}
G.~T.~Alcock,
``Gravity as Vacuum Electromagnetic Loading,''
DFD Research Programme (2026).

\bibitem{Touboul2022}
P.~Touboul \textit{et al.} (MICROSCOPE Collaboration),
``MICROSCOPE Mission: Final Results of the Test of the
Equivalence Principle,''
Phys.\ Rev.\ Lett.\ \textbf{129}, 121102 (2022).

\bibitem{Parker2018}
R.~H.~Parker, C.~Yu, W.~Zhong, B.~Estey, and H.~M\"uller,
``Measurement of the fine-structure constant as a test of
the Standard Model,''
Science \textbf{360}, 191 (2018).

\bibitem{Morel2020}
L.~Morel, Z.~Yao, P.~Clad\'e, and S.~Guellati-Kh\'elifa,
``Determination of the fine-structure constant with an accuracy
of 81 parts per trillion,''
Nature \textbf{588}, 61 (2020).

\bibitem{Fan2023}
X.~Fan, T.~G.~Myers, B.~A.~D.~Sukra, and G.~Gabrielse,
``Measurement of the Electron Magnetic Moment,''
Phys.\ Rev.\ Lett.\ \textbf{130}, 071801 (2023).

\bibitem{Hanneke2008}
D.~Hanneke, S.~Fogwell, and G.~Gabrielse,
``New Measurement of the Electron Magnetic Moment and the
Fine Structure Constant,''
Phys.\ Rev.\ Lett.\ \textbf{100}, 120801 (2008).

\end{thebibliography}

\end{document}
